% =========================================================================
% Remote Sensing of Environment (Elsevier) -- manuscript source (Phase F)
% Class: elsarticle (preprint, review, author-year citations, Harvard style)
% Grounding rule: every number must trace to docs/grounding/ (D1--D13).
% Unverified statements are marked with \citeneeded{...} or \pivalidate{...}.
% =========================================================================
\documentclass[preprint,review,12pt,authoryear]{elsarticle}

\usepackage[utf8]{inputenc}
\usepackage[T1]{fontenc}
\usepackage{amsmath,amssymb}
\usepackage{bm}
\usepackage{graphicx}
\usepackage{booktabs}
\usepackage{longtable}
\usepackage{multirow}
\usepackage{xcolor}
\usepackage[normalem]{ulem}
\usepackage{lineno}
\usepackage[hidelinks]{hyperref}

\graphicspath{{figures/}}
\linenumbers

% ---- Units and recurrent symbols ---------------------------------------
\newcommand{\um}{\,\ensuremath{\mu}m}
\newcommand{\gcm}{\,g\,cm\ensuremath{^{-2}}}
\newcommand{\NEdT}{NE\ensuremath{\Delta}T}
\newcommand{\Ldown}{L^{\downarrow}}
\newcommand{\Lup}{L^{\uparrow}}

% ---- Drafting markers (must be empty before submission) ----------------
\newcommand{\citeneeded}[1]{\textcolor{red}{[FIND REFERENCE: #1]}}
\newcommand{\pivalidate}[1]{\textcolor{blue}{[PI TO VALIDATE: #1]}}

\journal{Remote Sensing of Environment}

\begin{document}

\begin{frontmatter}

\title{Physics-informed neural networks with Sobolev shape regularization for thermal infrared temperature and emissivity separation: an open, leakage-controlled benchmark against the ASTER/HyTES TES algorithm}
% \pivalidate{Final title to be ratified after Results are drafted.}

\author[agro,ufrgs]{Atilio Grondona\corref{cor1}}
\ead{atilio.grondona@ufrgs.br}
\author[agro]{Fernanda Zunkowski Grondona}
\author[vlu]{Bijeesh Kozhikkodan Veettil}
\author[uchile]{Luis Morales-Salinas}
\cortext[cor1]{Corresponding author.}

% Author order and affiliations ratified by the PI on 12/11/2026 (4-author list synced with validated manuscript).
\affiliation[agro]{organization={Research and Development (R\&D) Department, AgroRisco AI},
            city={Porto Alegre}, state={RS}, postcode={91760-630}, country={Brazil}}
\affiliation[vlu]{organization={Department of Environmental Technology and Management, Faculty of Applied Technology, Van Lang University},
            city={Ho Chi Minh City}, postcode={700000}, country={Vietnam}}
\affiliation[uchile]{organization={Laboratory for Research in Environmental Sciences (LARES), Faculty of Agricultural Sciences, University of Chile},
            city={Santiago}, postcode={8820808}, country={Chile}}
\affiliation[ufrgs]{organization={Universidade Federal do Rio Grande do Sul (UFRGS)},
            city={Porto Alegre}, postcode={91501-970}, country={Brazil}}

\begin{abstract}
\textcolor{blue}{Retrieving Land Surface Temperature (LST) and Land Surface Emissivity (LSE) from thermal infrared radiances is a classically ill-posed inverse problem ($N$ channels for $N + 1$ physical unknowns). Operational algorithms, notably Temperature and Emissivity Separation (TES), resolve this Planck degeneracy using an empirical power law linking minimum emissivity to spectral contrast. While computationally expedient, this empirical heuristic amplifies channel noise into high-frequency sawtooth oscillations across hyperspectral channels, compounds atmospheric moisture errors, and degrades on anomalous mineralogies. We hypothesize that physics-informed neural networks regularized in function space by discrete Sobolev penalties and forward Planck consistency can resolve the Planck degeneracy without empirical contrast regressions. To test this hypothesis, we present a neural architecture featuring decoupled, thermodynamically bounded output heads ($T_s > 200\,\mathrm{K}$, $\varepsilon \in (0, 1)$), discrete Sobolev slope ($W^{1,1}$) and curvature ($W^{2,1}$) shape regularizers, and forward radiative consistency feedback. Using the open OpenTIR-PINN benchmark suite (240 curated laboratory spectra across 11 material classes, 30 MODTRAN~5.2 profiles, and a $2 \times 2$ double-quarantine evaluation matrix), we conduct a paired evaluation against operational TES across hyperspectral (HyTES, 223 channels) and multispectral (ASTER, 5 bands) regimes. In the hyperspectral regime, relying solely on surface-leaving radiance without auxiliary atmospheric telemetry, the network decisively outperforms classical TES: it reduces LST root-mean-square error by $44.0\%$ ($1.186$ vs.\ $2.119\,\mathrm{K}$), cuts channel-averaged emissivity error by $68.4\%$, halves water-vapor error propagation, and eliminates the 18-fold high-frequency noise amplification of classical TES to replicate natural spectral smoothness. Contiguous spectral sampling provides sufficient structural redundancy to constrain the physical inverse manifold directly. On multispectral ASTER, by contrast, sparse five-band sampling causes mathematical underdetermination that cannot decouple surface emission from atmospheric attenuation without auxiliary telemetry; here, the network achieves broad operational parity with classical TES in temperature while reconstructing emissivity shapes with superior angular fidelity. Stress-testing on quarantined golden targets demonstrates that on anomalous mineralogies exhibiting extreme absorption depths (exemplified by kaolinite), the neural prior induces cold temperature biases ($-4$ to $-12\,\mathrm{K}$). Model ensemble spread captures these out-of-domain targets on HyTES, but fails on ASTER, establishing the necessity of radiance-derived novelty filters before satellite deployment. Finally, batched inference executes tens of thousands of pixels per second on standard CPU hardware without GPU acceleration, demonstrating that physics-informed neural inversion is computationally tractable and ready for next-generation satellite missions, including NASA's SBG-TIR and ESA's LSTM.}
\end{abstract}

\begin{keyword}
Thermal infrared \sep Temperature and emissivity separation \sep Physics-informed neural networks \sep Sobolev regularization \sep HyTES \sep ASTER
\end{keyword}

\end{frontmatter}

\input{sections/01_introduction}
\input{sections/02_background}
\input{sections/03_dataset}
\input{sections/04_methods}
\input{sections/05_experimental_design}
\input{sections/06_results}
\input{sections/07_discussion}
\input{sections/08_conclusions}

\section*{Data and code availability}
To support community benchmarking, independent algorithmic validation, and further methodological development for thermal infrared missions, the complete OpenTIR-PINN v1.3.0 benchmark dataset, MODTRAN~5.2 radiative transfer look-up tables, 5-fold cross-validation split manifests, golden target spectra, pretrained model checkpoints, and reproduction scripts are publicly archived on Zenodo at \url{https://doi.org/10.5281/zenodo.24085242} (\url{https://zenodo.org/records/24085242}; concept DOI for all versions: \url{https://doi.org/10.5281/zenodo.23932830}) \citep{grondona2026opentir}. The open-source code repository, reproduction pipeline, 32-test validation suite, and interactive Streamlit analytical studio are available on GitHub at \url{https://github.com/atilioebg/OpenTIR-PINN}.

\section*{CRediT authorship contribution statement}
\pivalidate{CRediT roles for each author.}

\section*{Declaration of competing interest}
% Two authors are affiliated with a private company (AgroRisco AI). Elsevier requires disclosure of any
% financial interest that could be perceived as influencing the work.
\pivalidate{Confirm whether AgroRisco AI funded the work or has a commercial interest in TIR retrieval products; adapt the statement below accordingly.}
The authors declare that they have no known competing financial interests or personal relationships that could have appeared to influence the work reported in this paper.

\appendix
\input{sections/appA_modtran}
\input{sections/appB_validation}
\input{sections/appC_tes_sensitivity}

\bibliographystyle{elsarticle-harv}
\bibliography{references}

\end{document}

